% ══════════════════════════════════════════════════════════════════════════════
%  ADVANCED PARTICLE MODEL — the above-level book (from 2027)
%  Faculty decision, Nov 2026: every student completes the at-level lesson books;
%  students who excel move on to this book. Its topics go deeper into the ideas
%  of the at-level lessons rather than adding more of the same work.
%
%  Topics 1, 3, 4 and 6 are written here: a reading, the deeper idea worked
%  through (I do / We do), then %%PARTD:N%% pulls in the old Part D questions
%  of at-level Lesson N (both master lessons, where two make one) as practice,
%  then a Challenge. Topics 2 and 5 are the master's Lessons 6 and 12 (left out
%  of the at-level books), pulled in whole by %%LESSON:N%%.
%  Built by make_a5.py (via a5_levels.py).
% ══════════════════════════════════════════════════════════════════════════════

% ══════════════════════════════════════════════════════════════════════════════
%  TOPIC 1 — STRANGE STATES OF MATTER                        (deepens Lesson 1)
% ══════════════════════════════════════════════════════════════════════════════
\setcounter{section}{0}\renewcommand{\lessonart}{1}
\sectiondiv{Strange States of Matter}{Plasma \textbullet{} Glass and other amorphous solids \textbullet{} Fluids that change their minds \textbullet{} Deciding a state from evidence}\label{sec:1}
\lessoncodes{Deepens Lesson 1 \textbullet{} VC2S8U05 \textbullet{} VC2S10I06}

\vocabhead
\begin{ctab}
\begin{tabular}{|SW{3.4cm}|W{7.0cm}|}
  \hline\hd & \thead{Word} & \thead{Definition (copy from the slide)} \\ \hline
  & \scaf{\textbf{Plasma}} & \scaf{A gas so hot that its particles break apart into charged pieces; it glows and conducts electricity.} \\ \hline
  & \textbf{Amorphous solid} & \\ \hline
  & \textbf{Crystalline solid} & \\ \hline
  & \textbf{Viscosity} & \\ \hline
  & \textbf{Non-Newtonian fluid} & \\ \hline
\end{tabular}
\end{ctab}

\lessonpart{A}{Reading and comprehension}{Teal}\label{read:1}

\begin{readingbox}{Reading --- Not Quite Solid, Not Quite Liquid}
Solid, liquid and gas describe almost everything you meet. But the universe is stranger. The Sun and every other star are made of \textbf{plasma}: matter so hot that its particles have been torn into electrically charged pieces. Plasma glows and carries electricity. Most of the visible matter in the universe is plasma, which makes it, not solid, the most common state of all.

Even on Earth the three states have blurry edges. In salt, sugar and metals the particles are packed in neat, repeating rows: they are \textbf{crystalline solids}. Glass is different. Its particles are jumbled like a liquid's, but locked in place. Glass is an \textbf{amorphous solid}, a solid without an orderly pattern. It holds its shape and does not flow at room temperature, so it is a solid.

Liquids differ too. How thick or runny a liquid is, its resistance to flowing, is its \textbf{viscosity}. Honey has high viscosity; water has low viscosity. For most liquids the viscosity stays the same however hard you stir. But cornflour in water, called oobleck, gets thicker the harder you push it: it is a \textbf{non-Newtonian fluid}. Tomato sauce does the opposite and runs more easily once you shake it.
\end{readingbox}

\rcreset
\rcq{Why is plasma described as the most common state in the universe?}{2.5cm}
\rcq{How is glass different from salt, and why is it still called a solid?}{3cm}
\rcq{What is viscosity? Give one liquid with high and one with low viscosity.}{2.5cm}
\rcq{How do oobleck and tomato sauce behave differently when you push or shake them?}{2.5cm}

\lessonpart{B}{Theory and examples}{Blue}\label{theory:1}

\subsection{Deciding a State from Evidence}\label{sub:a1.1}

\begin{skillbox}{Four tests}
\textbf{Shape:} does it keep its own shape? \quad \textbf{Volume:} does it keep the same volume? \quad \textbf{Flow:} does it flow and take the shape of its container? \quad \textbf{Squash:} can it be compressed?\\
A solid passes the first two and fails the last two; a liquid keeps its volume, flows and cannot be squashed; a gas fails the first two, flows and can be squashed.
\end{skillbox}

\begin{workedex}{Honey}
Honey flows slowly into the shape of a jar, keeps its volume, and cannot be squashed.

\smallskip
It keeps its volume but not its shape, and it cannot be compressed, so honey is a \textbf{liquid}, one with a high viscosity. Flowing slowly does not make something a solid.
\end{workedex}

\begin{wedo}{Sand}
Dry sand can be poured from a bucket and takes the shape of the container. A student says sand must be a liquid. Use the four tests, and think about a single grain, to decide whether she is right.
\answerbox{5cm}
\end{wedo}

\subsection{Compressing a Gas, by Numbers}\label{sub:a1.2}

\begin{skillbox}{How many times smaller?}
\[ \text{times smaller} = \frac{\text{volume before}}{\text{volume after}} \]
Only gases can be squashed this much, because their particles are far apart with empty space between them.
\end{skillbox}

\begin{workedex}{A sealed syringe}
A sealed syringe holds 100\,mL of air. The plunger is pushed in until the air fills 25\,mL.

\smallskip
$100 \div 25 = 4$: the air is \textbf{4 times smaller}. The same number of particles now has a quarter of the space, so they are closer together and hit the walls more often.
\end{workedex}

\begin{wedo}{A gas bottle}
A barbecue gas bottle has a volume of 20\,L. The gas inside would fill about 4000\,L at normal air pressure. How many times smaller is it in the bottle? (In fact the gas is squashed so hard it turns into a liquid.)
\workbox{4cm}
\end{wedo}

\lessonpart{C}{Practice}{Amber}\label{prac:1}

%%PARTD:1%%

\lessonpart{D}{Challenge}{Violet}\label{ext:1}

\task A toy shop sells ``slime''. Plan an investigation to decide whether it behaves more like a solid or a liquid. Describe at least three tests, what you would do in each, and what result you would expect if it were a solid and if it were a liquid.
\answerbox{9cm}

\clearpage

% ══════════════════════════════════════════════════════════════════════════════
%  TOPIC 2 — MODELS AND THEIR LIMITS                     (deepens Lessons 2, 4)
%  The master's Lesson 6, whole.
% ══════════════════════════════════════════════════════════════════════════════
\setcounter{section}{1}\renewcommand{\lessonart}{6}
\sectiondiv{Models and Their Limits}{Common wrong ideas \textbullet{} Fixing a particle explanation \textbullet{} Checking a particle diagram \textbullet{} What the model cannot show}\label{sec:2}
\lessoncodes{Deepens Lessons 2 and 4 \textbullet{} VC2S8U05 \textbullet{} VC2S8H01}

%%LESSON:6%%

\clearpage

% ══════════════════════════════════════════════════════════════════════════════
%  TOPIC 3 — PARTICLES, ENERGY AND TEMPERATURE               (deepens Lesson 4)
% ══════════════════════════════════════════════════════════════════════════════
\setcounter{section}{2}\renewcommand{\lessonart}{4}
\sectiondiv{Particles, Energy and Temperature}{Proof that particles exist \textbullet{} Temperature and kinetic energy \textbullet{} Absolute zero and kelvin \textbullet{} How heat moves}\label{sec:3}
\lessoncodes{Deepens Lesson 4 \textbullet{} VC2S8U05 \textbullet{} VC2S10U14}

\vocabhead
\begin{ctab}
\begin{tabular}{|SW{3.4cm}|W{7.0cm}|}
  \hline\hd & \thead{Word} & \thead{Definition (copy from the slide)} \\ \hline
  & \scaf{\textbf{Kinetic energy}} & \scaf{The energy an object or particle has because it is moving.} \\ \hline
  & \textbf{Brownian motion} & \\ \hline
  & \textbf{Absolute zero} & \\ \hline
  & \textbf{Kelvin} & \\ \hline
  & \textbf{Thermal equilibrium} & \\ \hline
\end{tabular}
\end{ctab}

\lessonpart{A}{Reading and comprehension}{Teal}\label{read:3}

\begin{readingbox}{Reading --- The Dance of the Pollen Grains}
In 1827 the botanist Robert Brown watched tiny grains from inside pollen floating in water under his microscope. They jiggled about constantly, in random zigzags, and never stopped. This is now called \textbf{Brownian motion}. For almost eighty years nobody could explain it.

In 1905 Albert Einstein did. The grains, he said, are being struck from every side by water particles far too small to see. At any moment slightly more particles hit one side than the other, so the grain is shoved one way, then another. Einstein used the particle model to predict how far a grain should wander in a given time. In 1908 Jean Perrin measured it, and the numbers matched. Many scientists who had doubted that particles were real changed their minds.

Particles that move have \textbf{kinetic energy}. Temperature is a measure of the average kinetic energy of the particles: the hotter a substance, the faster its particles move. Cool something down and its particles slow. At \textbf{absolute zero}, about $-273$\,\textdegree C, they would have the least energy possible. Scientists measure from there using the \textbf{kelvin} scale, which starts at absolute zero.
\end{readingbox}

\rcreset
\rcq{What did Brown see, and why could he not explain it?}{2.5cm}
\rcq{Explain, in your own words, Einstein's explanation of Brownian motion.}{3cm}
\rcq{Why did Perrin's measurements change scientists' minds?}{2.5cm}
\rcq{What does temperature measure, in terms of particles?}{2cm}

\lessonpart{B}{Theory and examples}{Blue}\label{theory:3}

\subsection{The Kelvin Scale}\label{sub:a3.1}

\begin{skillbox}{Celsius and kelvin}
\[ \text{temperature in K} = \text{temperature in \textdegree C} + 273 \]
A change of 1\,\textdegree C is the same size as a change of 1\,K. There are no negative temperatures on the kelvin scale.
\end{skillbox}

\begin{workedex}{Converting}
Room temperature, 25\,\textdegree C: $25 + 273 = \mathbf{298\,K}$.\\
Liquid nitrogen boils at $-196$\,\textdegree C: $-196 + 273 = \mathbf{77\,K}$.
\end{workedex}

\begin{wedo}{Converting}
Convert to kelvin: water boiling (100\,\textdegree C); a freezer ($-18$\,\textdegree C). Convert to \textdegree C: 0\,K; the surface of the Sun, about 5800\,K.
\workbox{4.5cm}
\end{wedo}

\subsection{How Heat Moves}\label{sub:a3.2}

\begin{workedex}{A cold spoon in hot tea}
Explain, with particles, why a cold spoon left in hot tea warms up and the tea cools.

\smallskip
The tea's particles move fast and collide with the spoon's slower particles, making them vibrate faster. Each collision passes on some kinetic energy, so the spoon's particles speed up and the tea's slow down. Energy keeps moving from hot to cold until both are the same temperature. They are then in \textbf{thermal equilibrium}, and there is no more overall flow of energy.
\end{workedex}

\begin{wedo}{An ice cube in a drink}
Use particles to explain what happens to the ice cube and to the drink when an ice cube is dropped into a glass of water at room temperature. When does the change stop?
\answerbox{5cm}
\end{wedo}

\lessonpart{C}{Practice}{Amber}\label{prac:3}

%%PARTD:4%%

\lessonpart{D}{Challenge}{Violet}\label{ext:3}

\task A cup of water and a full bath are both at 40\,\textdegree C.
\pt{a} Compare the average kinetic energy of the particles in each.
\answerbox{3cm}
\pt{b} Which has more energy altogether? Explain.
\answerbox{3cm}
\pt{c} Which would melt more ice? Use your answers to explain why temperature and the total amount of thermal energy are not the same thing.
\answerbox{4cm}

\clearpage

% ══════════════════════════════════════════════════════════════════════════════
%  TOPIC 4 — ENERGY AND CHANGES OF STATE                     (deepens Lesson 5)
% ══════════════════════════════════════════════════════════════════════════════
\setcounter{section}{3}\renewcommand{\lessonart}{7}
\sectiondiv{Energy and Changes of State}{Why the temperature stops rising \textbullet{} Heating and cooling curves \textbullet{} Latent heat by numbers \textbullet{} Pressure and boiling}\label{sec:4}
\lessoncodes{Deepens Lesson 5 \textbullet{} VC2S8U05 \textbullet{} VC2S10U15}

\vocabhead
\begin{ctab}
\begin{tabular}{|SW{3.4cm}|W{7.0cm}|}
  \hline\hd & \thead{Word} & \thead{Definition (copy from the slide)} \\ \hline
  & \scaf{\textbf{Heating curve}} & \scaf{A graph of temperature against time as a substance is heated.} \\ \hline
  & \textbf{Latent heat} & \\ \hline
  & \textbf{Plateau} & \\ \hline
  & \textbf{Sublimation} & \\ \hline
  & \textbf{Pressure cooker} & \\ \hline
\end{tabular}
\end{ctab}

\lessonpart{A}{Reading and comprehension}{Teal}\label{read:4}

\begin{readingbox}{Reading --- Where Does the Energy Go?}
Put a pot of ice on a stove and measure its temperature every minute. The temperature climbs to 0\,\textdegree C and then stops, even though the stove is still heating it. It stays at 0\,\textdegree C until all the ice has melted. Then it climbs again, to 100\,\textdegree C, and stops again while the water boils. A graph of these readings is a \textbf{heating curve}, and the flat parts are \textbf{plateaus}.

Where does the energy go during a plateau? It is not making the particles move faster, so the temperature does not rise. Instead it is used to pull the particles apart against the attractions holding them together: from a fixed pattern into a liquid, or from a liquid into a gas. Energy used in this way is called \textbf{latent heat}, from a Latin word for ``hidden''.

Boiling point depends on the air pressure. High on a mountain, the air presses less on the water, so particles escape more easily and water boils below 100\,\textdegree C. A \textbf{pressure cooker} does the reverse: it traps steam, raises the pressure, and lets water reach about 120\,\textdegree C, so food cooks faster. Some solids skip the liquid state entirely: dry ice turns straight from solid to gas. This is \textbf{sublimation}.
\end{readingbox}

\rcreset
\rcq{Describe what happens to the temperature of ice as it is heated to boiling.}{2.5cm}
\rcq{During a plateau, what is the energy from the stove doing?}{2.5cm}
\rcq{Why does water boil below 100\,\textdegree C on a high mountain?}{2.5cm}
\rcq{How does a pressure cooker cook food faster?}{2.5cm}

\lessonpart{B}{Theory and examples}{Blue}\label{theory:4}

\subsection{Reading a Heating Curve}\label{sub:a4.1}

\begin{center}
\begin{tikzpicture}[font=\small, x=0.55cm, y=0.032cm]
  \draw[->, line width=0.7pt] (0,-30) -- (16,-30) node[below left, font=\scriptsize] {time (min)};
  \draw[->, line width=0.7pt] (0,-30) -- (0,130) node[above right, font=\scriptsize] {temperature (\textdegree C)};
  \foreach \y in {-20,0,50,100} \draw (-0.15,\y) -- (0,\y) node[left, font=\scriptsize] {\y};
  \draw[line width=1.2pt] (0,-20) -- (1.5,0) -- (5,0) -- (8,100) -- (13.5,100) -- (15,125);
  \node[font=\scriptsize] at (1.0,15) {P};
  \node[font=\scriptsize] at (3.3,10) {Q};
  \node[font=\scriptsize] at (6.0,55) {R};
  \node[font=\scriptsize] at (10.8,110) {S};
  \node[font=\scriptsize] at (14.0,125) {T};
\end{tikzpicture}
\figcap{Heating ice from $-20$\,\textdegree C until it has all boiled.}
\end{center}

\begin{workedex}{What is happening in each section?}
\textbf{P:} ice warming; particles vibrate faster. \quad \textbf{Q:} plateau at 0\,\textdegree C; ice melting, energy pulling particles out of their fixed pattern. \quad \textbf{R:} liquid water warming; particles move faster. \quad \textbf{S:} plateau at 100\,\textdegree C; water boiling, energy separating particles completely. \quad \textbf{T:} steam warming.
\end{workedex}

\begin{wedo}{A cooling curve}
Hot liquid wax cools from 90\,\textdegree C. Its temperature falls steadily, stays at 54\,\textdegree C for six minutes, then falls again to room temperature. Sketch the cooling curve, label the freezing point, and state what the particles are doing during the flat part. Is energy being taken in or given out there?
\drawbox{6cm}{}
\end{wedo}

\subsection{Latent Heat by Numbers}\label{sub:a4.2}

\begin{skillbox}{Energy to change state}
\[ \text{energy (J)} = \text{mass (g)} \times \text{latent heat (J/g)} \]
For water: melting needs 334\,J for every gram; boiling needs 2260\,J for every gram.
\end{skillbox}

\begin{workedex}{Melting and boiling 50\,g}
Melt 50\,g of ice: $50 \times 334 = 16\,700$\,J.\\
Boil away 50\,g of water: $50 \times 2260 = 113\,000$\,J, almost seven times as much.

\smallskip
Melting only loosens the particles; they still touch. Boiling must pull every particle completely away from its neighbours, which takes far more energy.
\end{workedex}

\begin{wedo}{An ice tray}
An ice tray holds 200\,g of water. How much energy must the freezer remove to freeze it once it is at 0\,\textdegree C? How much energy would boil the same water away?
\workbox{5cm}
\end{wedo}

\lessonpart{C}{Practice}{Amber}\label{prac:4}

%%PARTD:5%%

\lessonpart{D}{Challenge}{Violet}\label{ext:4}

\task A burn from steam at 100\,\textdegree C is much worse than a splash of boiling water at 100\,\textdegree C, even though they are at the same temperature. Explain why, using latent heat. Then calculate how much more energy 1\,g of steam gives your skin than 1\,g of boiling water, when each ends up as water at 100\,\textdegree C.
\answerbox{8cm}

\clearpage

% ══════════════════════════════════════════════════════════════════════════════
%  TOPIC 5 — GAS PRESSURE                                    (deepens Lesson 7)
%  The master's Lesson 12, whole.
% ══════════════════════════════════════════════════════════════════════════════
\setcounter{section}{4}\renewcommand{\lessonart}{12}
\sectiondiv{Gas Pressure}{Pressure from particle collisions \textbullet{} Squashing a gas \textbullet{} Heating a sealed gas \textbullet{} Air pressure all around us}\label{sec:5}
\lessoncodes{Deepens Lesson 7 \textbullet{} VC2S8U05 \textbullet{} VC2S8I01}

%%LESSON:12%%

\clearpage

% ══════════════════════════════════════════════════════════════════════════════
%  TOPIC 6 — DENSITY, FLOATING AND BUOYANCY                  (deepens Lesson 8)
% ══════════════════════════════════════════════════════════════════════════════
\setcounter{section}{5}\renewcommand{\lessonart}{13}
\sectiondiv{Density, Floating and Buoyancy}{Rearranging $D = m \div V$ \textbullet{} Average density \textbullet{} Archimedes' principle \textbullet{} Ships and submarines}\label{sec:6}
\lessoncodes{Deepens Lesson 8 \textbullet{} VC2S8U05 \textbullet{} VC2S10I04}

\vocabhead
\begin{ctab}
\begin{tabular}{|SW{3.4cm}|W{7.0cm}|}
  \hline\hd & \thead{Word} & \thead{Definition (copy from the slide)} \\ \hline
  & \scaf{\textbf{Buoyant force}} & \scaf{The upward push of a liquid or gas on an object in it.} \\ \hline
  & \textbf{Displace} & \\ \hline
  & \textbf{Average density} & \\ \hline
  & \textbf{Archimedes' principle} & \\ \hline
  & \textbf{Ballast} & \\ \hline
\end{tabular}
\end{ctab}

\lessonpart{A}{Reading and comprehension}{Teal}\label{read:6}

\begin{readingbox}{Reading --- Why a Steel Ship Floats}
Drop a steel bolt into water and it sinks. Yet a steel cargo ship weighing tens of thousands of tonnes floats. The secret is the shape. A ship's steel is spread into a thin hull around a huge space filled with air, so the \textbf{average density} of the whole ship, its total mass divided by its total volume, is less than the density of water.

When anything is placed in water, it pushes some water out of the way: it \textbf{displaces} it. The water pushes back up with a \textbf{buoyant force}. Over two thousand years ago, Archimedes worked out how big that force is. \textbf{Archimedes' principle} says the buoyant force equals the weight of the water displaced. A ship sinks into the water until it has displaced its own weight of water, then floats.

Sea water is denser than fresh water, so a ship floats a little higher at sea than in a river. A submarine controls whether it floats or sinks with \textbf{ballast} tanks: it lets sea water in to raise its average density and dive, and blows the water out with compressed air to rise.
\end{readingbox}

\rcreset
\rcq{Why does a steel ship float when a steel bolt sinks?}{2.5cm}
\rcq{State Archimedes' principle in your own words.}{2.5cm}
\rcq{Why does a ship float higher in sea water than in fresh water?}{2.5cm}
\rcq{How does a submarine change its average density to dive and rise?}{2.5cm}

\lessonpart{B}{Theory and examples}{Blue}\label{theory:6}

\subsection{Rearranging the Density Formula}\label{sub:a6.1}

\begin{skillbox}{Three forms of one rule}
\[ D = \frac{m}{V} \qquad m = D \times V \qquad V = \frac{m}{D} \]
Cover the quantity you want: what is left tells you whether to multiply or divide. Keep the units matched: g and cm\textsuperscript{3} give g/cm\textsuperscript{3}.
\end{skillbox}

\begin{workedex}{Finding a mass}
A copper block has a volume of 20\,cm\textsuperscript{3}. Copper has a density of 8.9\,g/cm\textsuperscript{3}. Find the mass.

\smallskip
$m = D \times V = 8.9 \times 20 = \mathbf{178\,g}$.
\end{workedex}

\begin{wedo}{Finding a volume}
A piece of cork has a mass of 12\,g. Cork has a density of 0.24\,g/cm\textsuperscript{3}. Find its volume.
\workbox{4cm}
\end{wedo}

\subsection{Will It Float?}\label{sub:a6.2}

\begin{workedex}{A sealed drum}
A sealed steel drum has a volume of 200\,000\,cm\textsuperscript{3} (200\,L) and a mass of 20\,kg (20\,000\,g). Does it float in fresh water (1\,g/cm\textsuperscript{3}), and how much cargo can it hold before it sinks?

\smallskip
Average density $= 20\,000 \div 200\,000 = \mathbf{0.1\ g/cm^3}$, less than water, so it floats.
It sinks when its average density reaches 1\,g/cm\textsuperscript{3}, that is, when its total mass reaches $1 \times 200\,000 = 200\,000$\,g. So it can carry $200\,000 - 20\,000 = 180\,000$\,g, or \textbf{180\,kg} of cargo.
\end{workedex}

\begin{wedo}{A canoe}
A canoe has a mass of 30\,kg and, if it were filled to the brim, would hold 300\,L (300\,000\,cm\textsuperscript{3}). What is the greatest total mass of paddlers and gear it could carry before water comes over the side? What happens to that answer in sea water (1.03\,g/cm\textsuperscript{3})?
\workbox{5.5cm}
\end{wedo}

\lessonpart{C}{Practice}{Amber}\label{prac:6}

%%PARTD:8%%

\needspace{10\baselineskip}
\task An egg sinks in a glass of tap water. As salt is stirred in, the egg slowly rises and then floats. Explain why, using density.
\answerbox{4.5cm}

\lessonpart{D}{Challenge}{Violet}\label{ext:6}

\task You need the density of a cork, which floats, so you cannot simply sink it in a measuring cylinder to find its volume. You also have a heavy metal sinker, string and a balance. Plan a method to find the cork's density, and show how you would calculate it from your measurements.
\answerbox{9cm}

\clearpage
